\documentclass[verbale=interno]{gpt}

\titolo{Primo incontro e costituzione del gruppo}
\data{2026-10-12}
\orario{21:00}{22:30}
\luogo{Discord}
\redattore{spagnoli}
\verificatore{corro}
\componenti{biasini,cescobolla,corro,guglielmo,mazzurana,spagnoli}
\assenti{biasini}

\modifica{1.0.0}{2026-10-14}{spagnoli}{corro}{Verifica e approvazione}
\modifica{0.1.0}{2026-10-13}{spagnoli}{--}{Prima stesura}

\begin{document}

\section{Ordine del giorno}
\begin{enumerate}
  \item Presentazione dei membri del gruppo.
  \item Scelta del nome del gruppo.
  \item Creazione dell'indirizzo email del gruppo.
  \item Scelta del logo.
  \item Strumenti di comunicazione del gruppo.
  \item Prima lettura dei capitolati.
  \item Organizzazione dell'incontro successivo.
\end{enumerate}

\section{Svolgimento}

\subsection{Presentazione dei membri}
L'incontro si è aperto con una presentazione dei membri presenti, per
fare conoscenza e iniziare a organizzare il lavoro comune. Enrico Biasini
non ha potuto partecipare.

\subsection{Nome, email e logo}
Il gruppo ha scelto il nome \textbf{Git Push Team} e ha creato
l'indirizzo \href{mailto:gitpushteam@gmail.com}{gitpushteam@gmail.com},
che sarà usato per tutte le comunicazioni ufficiali con docenti e
proponenti. È stato inoltre scelto il logo del gruppo, riportato in
Figura~\ref{fig:logo}.

\figurasenzabordo[0.28]{logo.png}{Logo del gruppo\label{fig:logo}}

\subsection{Strumenti di comunicazione}
Il gruppo ha stabilito di svolgere le riunioni su \textbf{Discord} e di
usare un gruppo \textbf{WhatsApp} per la messaggistica quotidiana.

\subsection{Ruoli}
In questo incontro non è stata ancora effettuata alcuna suddivisione dei
ruoli: il tema verrà affrontato nelle prossime riunioni.

\subsection{Prima lettura dei capitolati}
Il gruppo ha dato un primo sguardo ai capitolati presentati. Si è
stabilito che ciascun membro approfondisca individualmente i capitolati
che ritiene più interessanti, così da poterli presentare agli altri e
discuterne nell'incontro successivo.

\section{Decisioni}
\begin{decisioni}
  \decisione{Il nome del gruppo è \textbf{Git Push Team}.}
  \decisione{L'indirizzo email ufficiale del gruppo è
    gitpushteam@gmail.com.}
  \decisione{Viene adottato il logo riportato in Figura~\ref{fig:logo}.}
  \decisione{Le riunioni del gruppo si svolgono su Discord.}
  \decisione{Per la messaggistica il gruppo usa un gruppo WhatsApp.}
  \decisione{La suddivisione dei ruoli è rinviata a un incontro
    successivo.}
  \decisione{Ogni membro analizza individualmente i capitolati che lo
    interessano di più, in vista della discussione comune.}
\end{decisioni}

\section{Azioni da intraprendere}
\begin{azioni}
  \azione{Analizzare individualmente i capitolati di maggiore interesse
    e prepararsi a presentarli al gruppo.}{tutti}{13/10/2026}
\end{azioni}

\section{Prossima riunione}
Martedì 13 ottobre 2026 alle 17:40, per discutere i capitolati analizzati
da ciascun membro.

\end{document}
